\section{Introduction}
Cloud microservice systems have become the dominant architecture for modern online services \cite{11342355}, supporting large-scale applications such as e-commerce, finance, healthcare, and intelligent manufacturing \cite{10856683}. By decomposing close-coupled applications into loosely coupled services, microservice architectures significantly improve scalability, elasticity, and deployment efficiency. However, this architectural flexibility also introduces substantially higher operational complexity. A single user request may invoke dozens of dependent services, while failures can rapidly propagate across service chains and trigger cascading performance degradation\cite{zhang2025failure}. As a result, timely and accurate \textit{fault diagnosis} has become a fundamental requirement for ensuring service reliability and availability in cloud environments.

% Runtime fault diagnosis in cloud microservices is inherently a multimodal intelligence task. Numerical metrics reflect resource usage and performance fluctuations, logs record semantic evidence such as warnings and exceptions, while traces capture service dependencies and propagation paths. These heterogeneous signals provide complementary perspectives of system behavior, and effective diagnosis often requires jointly exploiting all of them. Motivated by this observation, recent studies have explored learning-based approaches for anomaly detection and root-cause analysis using multimodal operational data. Although promising progress has been achieved, existing methods still face important limitations when deployed in real-world latency-sensitive environments.

% First, advanced diagnosis models often emphasize predictive accuracy while incurring considerable inference overhead, making it difficult to support continuous online monitoring with low response latency. Second, cloud workloads and fault patterns are highly dynamic, whereas many existing methods rely on static architectures and uniform inference pipelines that cannot efficiently adapt computation to changing runtime contexts. Third, practical system operations require not only anomaly alarms, but also actionable diagnostic outputs such as probable root causes and interpretable evidence for rapid remediation. These limitations motivate the need for a new intelligent diagnosis framework that jointly considers effectiveness, efficiency, adaptability, and operational usability.

% To address the above challenges, this paper proposes Real-Time Multimodal Fault Diagnosis for Cloud Microservice Systems via Adaptive Sparse Inference, a practical AI-enabled diagnosis framework for modern cloud environments. The key idea is to integrate heterogeneous runtime observations while minimizing online inference cost through adaptive sparse computation. Specifically, we first design modality-specific encoders to model metrics, logs, and traces, and project them into a unified temporal representation space. A lightweight predictive backbone is then employed to capture normal system evolution patterns, where deviations between predicted and observed states are extracted as fault-sensitive representations. To improve deployment efficiency, we further develop an adaptive sparse inference mechanism that dynamically activates only a small subset of lightweight experts according to the current runtime context. Based on the learned representations, the framework jointly performs anomaly detection, root-cause ranking, and diagnosis-oriented interpretation.

% The main contributions of this paper are summarized as follows:

% A real-time multimodal diagnosis framework for cloud microservice systems.
% We develop a unified runtime diagnosis framework that jointly exploits metrics, logs, and traces for online anomaly detection and root-cause localization. By integrating heterogeneous operational signals into a single intelligent pipeline, the proposed method enables timely and accurate diagnosis for latency-sensitive cloud environments.
% An adaptive sparse inference mechanism for efficient AI deployment.
% We propose a resource-efficient inference architecture that combines a lightweight predictive backbone with context-aware sparse expert activation. The framework dynamically allocates computation according to runtime conditions, significantly reducing inference latency and computational overhead while maintaining strong diagnostic performance.
% A diagnosis-oriented intelligent operation capability with interpretable outputs.
% Beyond anomaly prediction, the proposed framework provides root-cause ranking and interpretable diagnostic evidence from multiple perspectives, improving the usability of AI-based fault diagnosis and supporting rapid operational decision-making in production systems.

% Extensive experiments on representative cloud microservice fault datasets demonstrate that the proposed framework consistently outperforms state-of-the-art baselines in fault detection accuracy and root-cause localization quality, while significantly reducing inference latency and computational overhead. These results validate the practicality and effectiveness of our design for real-world intelligent cloud operations.


Despite extensive progress in cloud observability, real-time fault diagnosis for microservice systems remains far from solved. Modern production platforms continuously generate massive volumes of heterogeneous monitoring signals, including metrics, logs, traces, and topology events. These data sources contain complementary information about system states, yet they differ significantly in temporal granularity, semantic structure, and noise characteristics, making unified diagnosis highly challenging \cite{10172617,brandon2020graph}. Meanwhile, diagnosis decisions are often required within seconds or even milliseconds after anomalies emerge, leaving little room for computationally expensive analysis pipelines. The challenge becomes more severe under dynamic workloads, where traffic bursts, service scaling, and evolving dependencies can quickly invalidate static diagnosis models and previously learned patterns \cite{pham2024baro,du2017deeplog}.

Existing approaches generally fall into three categories. Rule-based and statistical methods rely on handcrafted thresholds, invariant mining, or predefined causal patterns \cite{xu2009detecting,lou2010mining}, which are difficult to maintain in large and evolving systems. Supervised deep learning methods improve detection accuracy through representation learning \cite{du2017deeplog,meng2019loganomaly,10172617}, but typically require heavy end-to-end models whose inference latency grows rapidly with system scale. Recent foundation-model-based or large-model solutions provide stronger semantic understanding and reasoning capabilities \cite{guan2024logllm,wu2026icad}, yet their computational and deployment costs remain prohibitive for latency-sensitive online operations. More importantly, many existing methods optimize only diagnosis accuracy while overlooking real-time constraints such as bounded response delay, resource efficiency, and robustness under workload dynamics.

These limitations reveal three fundamental challenges for real-time cloud microservice diagnosis.

\textit{Challenge 1}: Heterogeneous multimodal system state modeling.
Metrics capture quantitative resource behavior, logs describe semantic execution events, and traces reveal cross-service invocation paths. How to fuse these heterogeneous signals into a unified and informative representation for online diagnosis remains difficult.

\textit{Challenge 2}: Low-latency diagnosis under resource constraints.
Online diagnosis engines must react quickly without competing excessively with production workloads for CPU, GPU, and memory resources. High diagnostic accuracy alone is insufficient if inference overhead is unacceptable.

\textit{Challenge 3}: Adaptation to dynamic environments.
Cloud systems continuously evolve due to autoscaling, workload shifts, software updates, and topology changes. Diagnosis models must remain effective under changing operating conditions rather than assuming static training and deployment distributions.

To address these challenges, we present ASID, an Adaptive Sparse Inference framework for real-time multimodal fault diagnosis in cloud microservice systems. The core idea of ASID is to combine expressive multimodal representation learning with runtime-efficient sparse computation. Specifically, ASID first employs modality-specific lightweight encoders to transform metrics, logs, and traces into a shared temporal representation space. It then learns predictive system dynamics and derives anomaly-sensitive residual features by measuring deviations between predicted and observed behaviors. Instead of activating a monolithic diagnosis network for every request, ASID introduces a context-aware sparse expert routing mechanism that dynamically selects only the most relevant diagnosis experts according to the current system state. This design preserves diagnostic capability while significantly reducing unnecessary computation. Based on the learned representations, ASID jointly performs anomaly detection, root-cause ranking, and interpretable diagnosis generation.

We implement ASID on a realistic cloud microservice diagnosis platform and evaluate it using representative fault scenarios under dynamic workloads. Results show that ASID consistently achieves superior diagnostic accuracy with substantially lower inference latency and computational overhead compared with state-of-the-art baselines, demonstrating its practicality for real-time intelligent cloud operations.

The main contributions of this paper are summarized as follows.
\begin{itemize}
\item We formulate the problem of real-time multimodal fault diagnosis for cloud microservice systems by jointly considering diagnostic accuracy, inference latency, and workload dynamics.
\item We propose ASID, a unified framework that integrates multimodal representation learning, predictive residual analysis, and context-aware sparse expert routing for efficient online diagnosis.
\item We develop an adaptive sparse inference mechanism that selectively activates diagnosis experts at runtime, significantly reducing computational overhead while maintaining high diagnostic quality.
\item We build a realistic evaluation platform with representative cloud fault scenarios and demonstrate that ASID outperforms existing methods in effectiveness, efficiency, and robustness.

\end{itemize}

The remainder of this paper is organized as follows. Section II reviews related work. Section III introduces the system model and problem formulation. Section IV presents the proposed framework in detail. Section V reports the experimental results and analysis. Finally, Section VI concludes this paper.